\documentclass[10pt,a4paper]{amsart}
\usepackage[margin=19mm]{geometry}
\usepackage{amsmath,amssymb,mathtools}
\usepackage[scr=rsfs,cal=euler]{mathalfa}
\usepackage{xfrac}
\usepackage[unicode,hidelinks,bookmarks=false]{hyperref}
\newcommand{\ie}{i.e.\ }
\newcommand{\cf}{cf.\ }
\newcommand{\resp}{respectively\ }
\newcommand{\Subsection}{\subsection}
\providecommand{\nobreakdash}{\nobreak\hbox{-}}
\setlength{\emergencystretch}{2em}
\multlinegap0pt
\usepackage{amssymb}
\usepackage{bm}
\usepackage{dsfont}
\usepackage{xcolor}
\usepackage{algorithm}\usepackage{algpseudocode}
\usepackage[normalem]{ulem}
\usepackage{cancel}
\usepackage{enumerate}
\usepackage[shortlabels]{enumitem}
\usepackage{tikz}
\newtheorem{theorem}{Theorem}
\newtheorem{lemma}{Lemma}
\usepackage{cleveref}
\crefname{theorem}{Theorem}{Theorems}
\crefname{lemma}{Lemma}{Lemmas}
\title{\textbf{Throughput-Optimal Scheduling Algorithms for LLM Inference and AI Agents}}
\author{J.G. Dai, Tianze Deng, Yueying Li, Tianyi Peng}
\date{Source: arXiv:2504.07347v3 (CC BY 4.0)}
\begin{document}
\maketitle

\noindent\emph{Source and licence.} \textbf{Throughput-Optimal Scheduling Algorithms for LLM Inference and AI Agents}. Version arXiv:2504.07347v3 (CC BY 4.0), distributed by arXiv under the Creative Commons Attribution 4.0 International licence, http://creativecommons.org/licenses/by/4.0/, as recorded in the arXiv licence metadata for this exact version and shown on the abstract page https://arxiv.org/abs/2504.07347v3. Full licence notice: \texttt{licenses/arxiv-2504.07347-CC-BY-4.0.txt}.\par\smallskip

\noindent\textit{Editorial scope.} Counted unit: the article's theorem convex\_hull (source0.tex lines 1220--1229). The article's complete original proof of it is delivered with it (source0.tex lines 1990--2068, one \texttt{proof} environment of 79 lines, which the article places away from the statement). No mathematics is written, repaired or re-derived: the statement and the proof are the article's own lines, in the article's own notation, and no retained line is substituted or re-flowed.  Retained uncounted as the source-local context the proof needs: lines 449--451; lines 1211--1214.  Nothing was omitted from any retained span, so the package carries no omission record.  No retained line contains a citation, so the wrapper carries no bibliography and no bibliography entry.  No retained line contains a figure environment, an image-inclusion command or drawing code, so no image is delivered and the package declares no asset dependency.  Editorial formatting only, in addition: an AMS \texttt{amsart} preamble that restates the article's own declarations and macros used by the retained lines, namely the environments \texttt{theorem}, \texttt{lemma} on separate counters, together with the standard packages those lines need.  The retained environments print in this document as Theorem 1, Lemma 1 in order of appearance, whereas the article prints the same items with the numbering of its own full document.  The complete native source of the article as distributed in the arXiv e-print for this exact version is bundled unchanged as \texttt{sources/176287/source0.tex}.
\begin{align}
    t_{b} = c + a \cdot \lceil \frac{b}{b_0} \rceil. \label{eq:batch-processing-time}  
\end{align}

\begin{align}
    \lambda_p &\le \sum_{\substack{x+y \le b_{\max} \\ y \le k_{\max}-1}} \frac{x}{t_{x+y}}\, u_{x,y}, \label{eq:stable_condition_lambda_p} \\
    \lambda_d &\le \sum_{\substack{x+y \le b_{\max} \\ y \le k_{\max}-1}} \frac{y}{t_{x+y}}\, u_{x,y} + \frac{k_{\max}}{t_{k_{\max}}}\, u_{0,k_{\max}}. \label{eq:stable_condition_lambda_d}
\end{align}

\begin{theorem} \label{convex_hull}
Assume batch processing times follow \cref{eq:batch-processing-time} with $a > 0$, $b_0 \mid b_{\max}$, and $k_{\max} \le b_0$. The pair $(\lambda_p, \lambda_d) \in \mathbb{R}^2_{\ge 0}$ satisfies the necessary condition \cref{eq:stable_condition_lambda_p}--\cref{eq:stable_condition_lambda_d} if and only if $(\lambda_p, \lambda_d)$ lies in the convex hull of the following four vertices:
\begin{align*}
    A &= \Bigl(\frac{b_{\max}}{t_{b_{\max}}},\, 0\Bigr), &
    B &= \Bigl(\frac{b_{\max}-k_{\max}+1}{t_{b_{\max}}},\, \frac{k_{\max}-1}{t_{b_{\max}}}\Bigr), \\
    C &= \Bigl(\frac{b_0-k_{\max}+1}{t_{b_0}},\, \frac{k_{\max}-1}{t_{b_0}}\Bigr), &
    D &= \Bigl(0,\, \frac{k_{\max}}{t_{k_{\max}}}\Bigr).
\end{align*}
Here, the inequality $(\lambda_p, \lambda_d) \le (a, b)$ is interpreted componentwise.
\end{theorem}

\begin{proof}[Proof of \cref{convex_hull}]
We prove the theorem in three steps: first establish a key monotonicity lemma, then verify the four vertices, and finally show that every feasible point lies in the convex hull.

\begin{lemma} \label{lem:monotonicity}
Let \( t_s = c + a \cdot \lceil s/b_0 \rceil \) with \( a > 0 \), \( c \ge 0 \), and \( b_0 > 0 \). Then for any integers \( x \ge 0 \), \( y \ge 0 \) with \( x + b_0 \le b_{\max} \),
\[
\frac{x + b_0}{t_{x + y + b_0}} \ge \frac{x}{t_{x + y}}.
\]
\end{lemma}

\begin{proof}
Let \( T = x + y \). Since \( \lceil (T + b_0)/b_0 \rceil = \lceil T/b_0 \rceil + 1 \), we have \( t_{T + b_0} = t_T + a \). Cross-multiplying, the desired inequality is equivalent to
\[
(x + b_0)\, t_T \ge x\, (t_T + a) \quad \Longleftrightarrow \quad b_0\, t_T \ge x\, a.
\]
Since \( \lceil T/b_0 \rceil \ge \lceil x/b_0 \rceil \ge x/b_0 \), we obtain
\[
b_0\, t_T = b_0\, c + a\, b_0 \lceil T/b_0 \rceil \ge b_0\, c + a\, x \ge a\, x,
\]
where the last inequality uses \( c \ge 0 \). This completes the proof.
\end{proof}

We now verify each vertex of the convex hull.

\paragraph{Vertex $A$.} The point \( A = \bigl( b_{\max}/t_{b_{\max}},\, 0 \bigr) \) corresponds to allocating all tokens to prefill. By \cref{lem:monotonicity} (with \( y = 0 \)), the ratio \( x/t_x \) satisfies \( (x + b_0)/t_{x+b_0} \ge x/t_x \) for all \( x \le b_{\max} - b_0 \). Applying this inductively from \( x = 0 \) in steps of \( b_0 \), and using \( b_0 \mid b_{\max} \), we conclude that \( b_{\max}/t_{b_{\max}} \) is the maximum achievable prefill rate. Hence \( A \) is a vertex.

\paragraph{Vertex $D$.} The point \( D = \bigl( 0,\, k_{\max}/t_{k_{\max}} \bigr) \) corresponds to decode-only batches using all \( k_{\max} \) slots. Since \( k_{\max} \le b_0 \), we have \( \lceil y/b_0 \rceil = 1 \) for all \( 1 \le y \le k_{\max} \), so \( t_y = c + a \) is constant in this range. Thus \( y/t_y = y/(c + a) \) is increasing in \( y \), and the decode rate is maximized at \( y = k_{\max} \). Hence \( D \) is a vertex.

\paragraph{Vertices $B$ and $C$.} These represent mixed batches with \( y = k_{\max} - 1 \) decode tokens at different prefill loads:
\begin{align*}
    B &= \Bigl(\frac{b_{\max} - k_{\max} + 1}{t_{b_{\max}}},\, \frac{k_{\max} - 1}{t_{b_{\max}}}\Bigr), &
    C &= \Bigl(\frac{b_0 - k_{\max} + 1}{t_{b_0}},\, \frac{k_{\max} - 1}{t_{b_0}}\Bigr).
\end{align*}

\paragraph{Every feasible point lies in the convex hull.}
Consider any feasible operating point \( (x/t_{x+y},\, y/t_{x+y}) \) with \( x + y \le b_{\max} \) and \( y \le k_{\max} - 1 \). Let \( m \) be the unique integer such that \( (m-1)b_0 < x + y \le m b_0 \le b_{\max} \). Then \( t_{x+y} = t_{m b_0} = c + am \), and since \( x \le m b_0 - y \),
\[
\frac{x}{t_{x+y}} \le \frac{m b_0 - y}{t_{m b_0}}, \qquad \frac{y}{t_{x+y}} \le \frac{y}{t_{m b_0}}.
\]
Thus the point is dominated (componentwise) by
\[
P_m = \Bigl( \frac{m b_0 - y}{t_{m b_0}},\, \frac{y}{t_{m b_0}} \Bigr).
\]
It suffices to show \( P_m \in \mathrm{conv}(O, A, B, C, D) \) for all valid \( m \) and \( 0 \le y \le k_{\max} - 1 \).

\textbf{Step 1.} For each fixed \( m \), define
\[
A_m = \Bigl( \frac{m b_0}{t_{m b_0}},\, 0 \Bigr), \qquad Q_m = \Bigl( \frac{m b_0 - k_{\max} + 1}{t_{m b_0}},\, \frac{k_{\max} - 1}{t_{m b_0}} \Bigr).
\]
Then \( P_m \) is a convex combination of \( A_m \) and \( Q_m \):
\[
P_m = \Bigl(1 - \frac{y}{k_{\max} - 1}\Bigr) A_m + \frac{y}{k_{\max} - 1}\, Q_m.
\]
By \cref{lem:monotonicity}, the sequence \( m b_0 / t_{m b_0} \) is nondecreasing in \( m \), so \( A_m \) lies on the segment \( \overline{OA} \) and hence in the convex hull. It remains to show \( Q_m \) lies in the convex hull.

\textbf{Step 2.} We show that \( Q_m \) lies on or below the line segment \( \overline{BC} \) for all \( 1 \le m \le b_{\max}/b_0 \). Note that \( Q_1 = C \) and \( Q_{b_{\max}/b_0} = B \), so the boundary cases hold with equality. For \( 1 < m < b_{\max}/b_0 \), let \( q = b_{\max}/b_0 \). The \( x \)-coordinates of \( B \), \( C \), and \( Q_m \) are
\[
x_B = \frac{q b_0 - k_{\max} + 1}{c + aq}, \quad x_C = \frac{b_0 - k_{\max} + 1}{c + a}, \quad x_{Q_m} = \frac{m b_0 - k_{\max} + 1}{c + am}.
\]
By \cref{lem:monotonicity}, \( x_{Q_m} \) is increasing in \( m \), so \( x_C \le x_{Q_m} \le x_B \). Thus there exists a unique \( \alpha \in [0,1] \) with \( x_{Q_m} = \alpha\, x_B + (1 - \alpha)\, x_C \).

The \( y \)-coordinate of \( Q_m \) is \( (k_{\max} - 1)/(c + am) \), while the corresponding point on \( \overline{BC} \) has \( y \)-coordinate
\[
y_{\mathrm{line}} = (k_{\max} - 1) \Bigl( \frac{\alpha}{c + aq} + \frac{1 - \alpha}{c + a} \Bigr).
\]
Define \( f(m) = 1/(c + am) \). Since \( a > 0 \), we have \( f''(m) = 2a^2/(c + am)^3 > 0 \), so \( f \) is strictly convex. By Jensen's inequality,
\[
f(m) < \alpha\, f(q) + (1 - \alpha)\, f(1)
\]
for \( 1 < m < q \), which gives
\[
\frac{k_{\max} - 1}{c + am} < (k_{\max} - 1) \Bigl( \frac{\alpha}{c + aq} + \frac{1 - \alpha}{c + a} \Bigr).
\]
Hence the \( y \)-coordinate of \( Q_m \) is strictly below \( y_{\mathrm{line}} \), confirming \( Q_m \) lies strictly below \( \overline{BC} \) and therefore inside \( \mathrm{conv}(O, A, B, C, D) \).

\textbf{Step 3: Decode-only batches.} Finally, consider decode-only batches with \( x = 0 \) and \( y = k_{\max} \). The operating point is \( (0, k_{\max}/t_{k_{\max}}) = D \), which is a vertex of the convex hull.

\textbf{Conclusion.} Every feasible point \( (x/t_{x+y},\, y/t_{x+y}) \) with \( x + y \le b_{\max} \) and \( y \le k_{\max} \) is dominated by a point in \( \mathrm{conv}(O, A, B, C, D) \). Conversely, each vertex is achievable by an appropriate batch composition, so the convex hull is tight.
\end{proof}

\end{document}
